% ============================================================================
%  SECCIÓN 3.14: ESTRUCTURA INICIAL DEL PROYECTO
% ============================================================================

\section{Estructura del proyecto}

La estructura del proyecto define la organizaci\'on de los archivos y
directorios que componen la aplicaci\'on, facilitando su mantenimiento,
escalabilidad y la incorporaci\'on de nuevas funcionalidades. Para el sistema
propuesto se adopt\'o una organizaci\'on profesional por capas, en la que el
c\'odigo fuente se agrupa en el directorio \texttt{src/} y cada responsabilidad
(caracter\'isticas, servicios, n\'ucleo) se separa de forma clara.

La estructura inicial del proyecto es la siguiente:

\begin{lstlisting}
lugares-interactivos/
|-- assets/
|   |-- icon.png
|   |-- favicon.png
|   |-- android-icon-background.png
|   |-- android-icon-foreground.png
|   `-- android-icon-monochrome.png
|-- src/
|   |-- app/
|   |   |-- navigation/
|   |   |   |-- AppNavigator.tsx
|   |   |   |-- RootNavigator.tsx
|   |   |   `-- types.ts
|   |   |-- screens/
|   |   |   |-- auth/
|   |   |   |   |-- LoginScreen.tsx
|   |   |   |   `-- RegisterScreen.tsx
|   |   |   |-- home/
|   |   |   |   |-- HomeScreen.tsx
|   |   |   |   `-- components/
|   |   |   |       `-- ScenarioCard.tsx
|   |   |   |-- map/
|   |   |   |   |-- MapScreen.tsx
|   |   |   |   `-- components/
|   |   |   |       `-- ScenarioMap.tsx
|   |   |   |-- catalog/
|   |   |   |   |-- CatalogScreen.tsx
|   |   |   |   `-- components/
|   |   |   |       `-- FilterBar.tsx
|   |   |   `-- details/
|   |   |       |-- ScenarioDetailScreen.tsx
|   |   |       `-- components/
|   |   |           `-- InfoSection.tsx
|   |   |-- context/
|   |   |   |-- AuthContext.tsx
|   |   |   |-- FavoritesContext.tsx
|   |   |   `-- ThemeContext.tsx
|   |   |-- components/
|   |   |   |-- common/
|   |   |   |   |-- Button.tsx
|   |   |   |   |-- Input.tsx
|   |   |   |   `-- LoadingIndicator.tsx
|   |   |   `-- layout/
|   |   |       `-- Header.tsx
|   |   |-- services/
|   |   |   |-- api/
|   |   |   |   |-- client.ts
|   |   |   |   |-- scenarios.ts
|   |   |   |   `-- auth.ts
|   |   |   |-- storage/
|   |   |   |   `-- tokenStorage.ts
|   |   |   `-- location/
|   |   |       `-- locationService.ts
|   |   |-- hooks/
|   |   |   |-- useScenarios.ts
|   |   |   |-- useFavorites.ts
|   |   |   `-- useLocation.ts
|   |   |-- types/
|   |   |   |-- scenario.ts
|   |   |   |-- user.ts
|   |   |   `-- navigation.ts
|   |   `-- theme/
|   |       |-- colors.ts
|   |       |-- spacing.ts
|   |       `-- index.ts
|   |-- App.tsx
|   |-- index.ts
|-- .env.example
|-- .gitignore
|-- app.json
|-- package.json
|-- pnpm-lock.yaml
|-- tsconfig.json
`-- README.md
\end{lstlisting}

Esta organizaci\'on sigue el principio de separaci\'on de responsabilidades, de
modo que cada directorio cumple una funci\'on espec\'ifica y los componentes se
reutilizan entre las distintas pantallas \citep{sommerville2011}. En la tabla
se describen los directorios principales de la estructura del proyecto.

\begin{table}[H]
  \centering
  \caption{Directorios principales de la estructura del proyecto}
  \contenidoConFuente{%
    \begin{tablaAPA}
      \begin{tabular}{@{}p{3.2cm}p{12.6cm}@{}}
        \toprule
        \textbf{Directorio} & \textbf{Funci\'on} \\
        \midrule
        \texttt{assets/} &
        Recursos est\'aticos de la aplicaci\'on: iconos, im\'agenes y logotipos. \\
        \addlinespace
        \texttt{src/app/navigation/} &
        Configuraci\'on de la navegaci\'on de la aplicaci\'on: pilas y pesta\~nas,
        definiendo las rutas entre las pantallas. \\
        \addlinespace
        \texttt{src/app/screens/} &
        Pantallas de la aplicaci\'on, agrupadas por m\'odulo funcional: autenticaci\'on,
        inicio, mapa, cat\'alogo y detalle. \\
        \addlinespace
        \texttt{src/app/context/} &
        Contextos globales de React para el estado compartido: autenticaci\'on,
        favoritos y tema. \\
        \addlinespace
        \texttt{src/app/components/} &
        Componentes reutilizables de interfaz, separados en elementos comunes y de
        dise\~no (botones, campos, encabezados). \\
        \addlinespace
        \texttt{src/app/services/} &
        Servicios de comunicaci\'on con el backend: cliente de API, autenticaci\'on,
        almacenamiento local y servicios de ubicaci\'on. \\
        \addlinespace
        \texttt{src/app/hooks/} &
        Hooks personalizados de React que encapsulan la l\'ogica reutilizable de
        datos y ubicaci\'on. \\
        \addlinespace
        \texttt{src/app/types/} &
        Definici\'on de los tipos y modelos de datos de la aplicaci\'on
        (escenarios, usuarios, navegaci\'on). \\
        \addlinespace
        \texttt{src/app/theme/} &
        Tema visual de la aplicaci\'on: colores, espacios y estilos comunes. \\
        \addlinespace
        \texttt{App.tsx} &
        Componente ra\'iz de la aplicaci\'on que integra los proveedores de contexto
        y la navegaci\'on. \\
        \addlinespace
        \texttt{index.ts} &
        Punto de entrada de la aplicaci\'on que registra el componente ra\'iz en
        React Native. \\
        \addlinespace
        \texttt{app.json} &
        Archivo de configuraci\'on de Expo: nombre, identificador, orientaci\'on,
        iconos y ajustes de la plataforma \citep{expo}. \\
        \addlinespace
        \texttt{package.json} &
        Declara las dependencias del proyecto y los scripts de ejecuci\'on. \\
        \addlinespace
        \texttt{.env.example} &
        Plantilla de las variables de entorno necesarias para la configuraci\'on
        de los servicios externos. \\
        \bottomrule
      \end{tabular}
    \end{tablaAPA}%
  }{Elaboraci\'on propia en base al repositorio del proyecto.}
\end{table}

La organizaci\'on modular por capas permite que cada m\'odulo del sistema (mapa
interactivo, cat\'alogo, autenticaci\'on, favoritos) se desarrolle, pruebe y
mantenga de forma independiente, tal como lo exigen los requerimientos no
funcionales de portabilidad y mantenibilidad del sistema \citep{iso25010}.
